\section{Experimental Setup}
\label{sec:experimental-setup}

This section defines the evaluation protocol used to support the paper's claims. The experiments are organized around one main scenario and a small set of supporting studies. The main scenario is the dual-formation eight-sensor distributed GA-LMB tracking problem under tiered heterogeneous packet loss. Supporting studies check whether the observed gains persist under ideal communication, how the method behaves as communication becomes more constrained, and how optional modules compare with the retained design.

\subsection{Scenario, protocols, and operating points}

The evaluation follows a hierarchy. The main scenario uses an eight-sensor formation with tiered packet loss and provides the confirmatory evidence. A matched factor ablation identifies which parts of the weight design carry the result. Runtime in the same scenario measures deployment cost. An ideal-communication study tests whether the refinements remain useful without packet loss. A four-level sensitivity study then examines Balanced mode as communication becomes more constrained. The appendix reports exploratory variants.

\paragraph{Main scenario.}

The primary evaluation setting uses eight mobile sensors arranged as two four-sensor formations. The filtering mode is distributed local fusion. Each node forms a neighborhood problem from its own and its neighbors' delivered measurement streams. It then obtains measurement-updated LMB posteriors for those streams and fuses the resulting parameters. The local multi-object filter is GA-LMB with loopy belief propagation (LBP) for approximate marginal data association \cite{WilliamsLau2014LBP}.

The choice of the geometric-average KLA route is deliberate. GA fusion gives a conservative logarithmic-opinion-pool update under unknown cross-correlations. It also admits a stable canonical-form update for LMB spatial densities and exposes the Bernoulli product rule needed for branch-specific weighting. These properties make it a cleaner vehicle for studying adaptive weight allocation than a broad comparison across fusion families.

The common sensor-side configuration is:
\begin{itemize}
\item number of sensors: $8$,
\item filter mode: GA-LMB,
\item data association: LBP,
\item clutter rate: $3$ for every sensor,
\item detection probability: $0.9$ for every sensor,
\item measurement-noise scale: $q = 3$ for every sensor,
\item sensor field of view enabled with half-angle $60^\circ$ and range $60000$.
\end{itemize}
The nominal detection probability, measurement-noise scale, and clutter rate are therefore homogeneous across sensors. The main sensing variation comes from the dual-formation geometry, finite fields of view, and target visibility. These conditions induce differences in posterior covariance and existence, but they do not vary explicit $p_D$, measurement-noise, or clutter tiers.

The target side consists of ten objects arranged in three groups with staggered births. The three groups contain $3 + 3 + 4$ targets, respectively, and the targets move toward the central region with nominal speed $0.45$. Births are staggered every eight time steps, and the total simulation length is $100$ steps.

\paragraph{Communication model.}

The main scenario uses a tiered heterogeneous packet-loss setting. Communication is applied before distributed fusion through a measurement-stream delivery model with bandwidth limitation and per-sensor packet loss. The main configuration is:
\begin{itemize}
\item communication level: $2$,
\item global maximum delivered measurements per step: $80$,
\item packet-selection policy: weighted priority with random measurement selection,
\item link model: fixed packet drop,
\item target mean drop rate: $p_{\mathrm{drop}} = 0.2$,
\item tiered drop levels: $[0, 0.1, 0.2, 0.5]$,
\item tiered level counts: $[1, 4, 1, 2]$.
\end{itemize}
This tiered design preserves the historical mean communication degradation while introducing persistent cross-node heterogeneity.

Figure~\ref{fig:scenario-schematic} summarizes the contrast between the simpler standard fixed ideal benchmark and the main tiered-loss scenario. The standard ideal benchmark uses four full region-of-interest (ROI) sensors, ideal communication, and fixed target births. The main scenario is the dual-formation eight-sensor case with staggered target waves, a bandwidth cap, and persistent drop-tier heterogeneity before distributed fusion.

\begin{figure}[t]
  \centering
  \includegraphics[width=\linewidth]{paper-figure-scenario.pdf}
  \caption{Schematic comparison between the standard fixed ideal benchmark and the main tiered heterogeneous packet-loss scenario. The ideal benchmark uses four full-ROI sensors with no packet loss and fixed target births; the main scenario uses the dual-formation eight-sensor geometry and staggered target waves with the bandwidth cap and persistent drop-tier heterogeneity used in the main experiments.}
  \label{fig:scenario-schematic}
\end{figure}

\paragraph{Mainline ablation protocol.}

The core ablation study compares five consistently named arms:
\begin{enumerate}
\item Fixed Metropolis,
\item Covariance-only adaptive,
\item Covariance-link adaptive,
\item Three-factor adaptive backbone,
\item Balanced mode, which adds weak structure-aware modulation after covariance, link-quality, and existence-confidence weighting.
\end{enumerate}

This five-arm path diagnoses the backbone before the final existence refinement. The results table focuses on operating points and treats the three-factor backbone as part of Balanced mode. The main comparison also includes probability-of-detection (PD)-weighted GA, scalar FID-FIA-weighted GA, and Cardinality-critical mode. Balanced mode uses one parameterization across all trials. Exponential smoothing is $\alpha_{\mathrm{EMA}}=0.7$, and the minimum weight is $0.05$. Existence-confidence parameters are $\lambda_{\min}=0.85$ and $p_e=2.0$. Spatial and existence interpolation strengths are $0.5$ and $0.15$, and the graph-prior strengths are $0.45$ and $0.08$. The nominal communication prior uses power $0.30$ and lower bound $0.25$. Cardinality-critical mode leaves the spatial branch unchanged but adds FID-FIA to existence. Smoothing and minimum-weight flooring are disabled for this final existence refinement. Other exploratory modules are inactive.

\paragraph{Reviewer-requested external-baseline adaptations.}
Two additional rows map the closest peer-to-peer methods onto the same one-round GA-LMB protocol \cite{Zheng2026DistributedWeightedRFS,Gao2023ResilientLMB}. We label both rows as adaptations because their source methods use a different filter implementation or threat model. Zheng et al. formulate weighted RFS subdensity fusion and provide consensus and flooding implementations for GM-PHD filters \cite{Zheng2026DistributedWeightedRFS}. They do not provide an LMB implementation. In our adaptation, each label-associated Bernoulli component is one subdensity. For label $\ell$ and active neighbor $j$, the spatial and existence coefficients are
\begin{equation}
 a_{x,j}^{(\ell)}=\frac{1}{\operatorname{tr}(P_{j,\mathrm{pos}}^{(\ell)})+\epsilon},
 \qquad
 a_{r,j}^{(\ell)}=\epsilon_r+(1-\epsilon_r)\left|2r_j^{(\ell)}-1\right|,
 \label{eq:zheng-style-adaptation}
\end{equation}
where $\epsilon$ is machine precision and $\epsilon_r=0.05$. Each coefficient vector is normalized over available neighbors and applied to the corresponding GA-LMB branch. The adaptation retains subdensity-level weighting and separate localization and cardinality cues. It is not equivalent to the source GM-PHD feedback or flooding implementations.

Gao et al. address corrupted LMB messages and adversarial nodes \cite{Gao2023ResilientLMB}. Each receiving node treats its local density as trusted, detects modified or false Bernoulli components, and reduces confidence as discrepancy grows. Our adaptation aligns components by label and assigns neighbor $j$ the credibility
\begin{equation}
 c_{s,j}^{(\ell)}=\max\!\left\{\epsilon_c,
 \exp\!\left[-\frac{J_{\mathrm{G},s,j}^{(\ell)}/d_x+J_{\mathrm{B},s,j}^{(\ell)}}{\tau}\right]\right\},
 \qquad
 \omega_{s,j}^{(\ell)}\propto \bar{\omega}_{s,j}c_{s,j}^{(\ell)},
 \label{eq:gao-style-adaptation}
\end{equation}
Here, $J_{\mathrm{G}}$ and $J_{\mathrm{B}}$ are symmetric KL divergences for the projected Gaussian and Bernoulli distributions. We set $d_x=4$, $\tau=2$, and $\epsilon_c=0.05$. The coefficient $\bar{\omega}_{s,j}$ is the Metropolis weight, and the receiving node has unit credibility. Hard rejection is disabled because the experiment modifies packet delivery but not posterior parameters. This row therefore evaluates soft local-trust weighting, not the full attack detector. Every adaptation and proposed method uses the same truth, measurements, delivery realization, and 50 deterministic paired seeds $2$--$51$. The adaptation arms and the reused core arms are executed in separate run cohorts. Their accuracy metrics remain paired by seed, but we do not report cross-cohort relative runtime.

\subsection{Metrics and reporting}

The experiments report three categories of metrics. The first category is network-level disagreement among sensors after distributed fusion. These primary metrics are OSPA consensus error, GOSPA consensus error, matched localization disagreement, and cardinality dispersion as defined in Section~\ref{sec:problem-formulation}. The second category is conventional truth-referenced tracking accuracy. These secondary metrics are local Euclidean OSPA (E-OSPA), local root mean square error (RMSE), and local cardinality error. The third category is computational cost, measured as filter/fusion wall-clock runtime per trial, runtime per simulation step, and runtime relative to Fixed Metropolis. Every Euclidean OSPA and GOSPA calculation uses cutoff $c=5$ and order $p=2$, with $\alpha=2$ for GOSPA. Their ground distance uses the same complete extracted kinematic state vectors. OSPA and GOSPA provide complementary normalized and unnormalized finite-set distances for multi-object filtering and tracking evaluation \cite{Schuhmacher2008OSPA,Rahmathullah2017GOSPA,Ristic2011TrackMetric}. The primary disagreement metrics are task-specific network-agreement measures built from these standard ingredients, and the paper interprets them together with the secondary local metrics and computational-cost measurements.

This hierarchy reflects the method's purpose. The method is not a generic single-sensor accuracy booster. It reduces disagreement in the fused multisensor picture under heterogeneous communication. Network disagreement is therefore the main outcome. Local tracking metrics guard against hidden accuracy losses, and runtime measures deployment cost. The target is better agreement without a meaningful loss in conventional tracking quality or impractical computation. Runtime covers only the distributed LMB filtering and fusion call for each arm. Scenario generation, communication sampling, plotting, and metric evaluation are excluded.

Unless otherwise stated, reported scalar metrics are Monte Carlo means. The main fixed-versus-existence-refined confirmation, the five-arm factor ablation, the ideal-communication comparison, and the communication-level sensitivity study all use 50 deterministic paired trials with seeds $2$--$51$. The main confirmation, ideal-communication comparison, and communication-level sensitivity study are paired by construction. The five-arm factor ablation remains primarily a design diagnostic, but it is now reported on the same 50-trial seed set for consistency.
